\Chapter{Defining Natural Numbers}
\Lettrine[loversize=0.1,nindent=0pt]{To} define all the numbers, we have to start with the most basic of numbers -- the natural numbers. The natural numbers are the most basic of all numbers. They are considered our counting numbers, and understanding how they can be defined using the axioms of Zermelo-Frankel theory is essential for continuing to create the rest of the numbers.

\section{The Need for Abstraction}
So why go through all this trouble to define the natural numbers? Surely everyone can agree what a natural number is?

Let us ask a simple question: what is 3? Seriously. Think about what 3 really means. In elementary school, we were taught that 3 is just a number quantifying a certain number of things. There are many 3s -- 3 circles, 3 squares, 3 cars, 3 stars, 3 apples, 3 lines, 3 words, 3 points, 3 boxes, etc.
\begin{figure}[H]
    \centering
    \includegraphics[width=8cm]{part-1/constructing-numbers/three.png}
    \caption{Various Examples of 3}
\end{figure}
What they all encapsulate is that there is a certain number of things. And \textbf{that} number is 3. It is an abstraction above a specific `type' of 3. Not 3 circles, or 3 squares, or any of the other types that 3 could be; 3 is a pure symbol symbolising a pure count.

Of course, when it comes to asking what 3 is, at a fundamental level, we need to have a representative for the concept of 3. This is just like how Physics moved away from defining the SI base units using physical objects and their properties, and instead defining them using fundamental physical constants of the Universe. Although they have shifted the burden away from a specific instance of the physical artefact, they still need one for practical use. Likewise, although we can abstract away the idea of 3, we still need a representative for 3.

In this chapter, we will be focusing on defining the natural numbers using a set-theoretic approach. In particular, we will use \ZF~to construct representatives for \textbf{all} the natural numbers. That way, we will have no ambiguity when we discuss what 3, or 12, or 83, or 1289, or 89338, etc., really means at a fundamental level.

\section{Defining Zero}
TODO: Add

\section{Successor}
TODO: Add

\section{The Natural Numbers}
TODO: Add

\section{Axiom of Infinity}
TODO: Add
% Probably need to motivate the reason for how this axiom is defined
